\vfill\pagebreak
\section{Several functions with the same name}
\label{sec:functions:overloading}

The area of a square needs one length, and the area of a rectangle two. Both
functions compute an area, and names such as \token{squareArea} and
\token{rectangleArea} only repeat what the arguments already say. Ymir lets
several functions share a name, as long as their parameters differ, in number or
in type. The name is then \textit{overloaded}, and each call is sent to the
function whose parameters match its arguments.

\lstinputlisting[style=coloredverbatim, caption={Two functions named \token{area}, \textit{examples/functions/area.yr}}, label=lst:functions:area]{examples/functions/area.yr}
\vspace{-10pt}%
\begin{lstlisting}[style=bashVerb]
A 3 by 3 square: 9
A 2 by 5 rectangle: 10
\end{lstlisting}

The call on line~12 has one argument, and only the first \token{area} takes one.
The call on line~13 has two, and goes to the second. The choice is made by the
compiler, once, when it reads the call, not each time the program runs.

Functions with the same number of parameters are told apart by their types. The
three functions of the next listing take one argument each, of three different
types, and each call goes to the function whose parameter has the type of its
argument.

\lstinputlisting[style=coloredverbatim, caption={Overloading on the type of the argument, \textit{examples/functions/describe.yr}}, label=lst:functions:describe]{examples/functions/describe.yr}
\vspace{-10pt}%
\begin{lstlisting}[style=bashVerb]
42 is an integer
2.5 is a floating point number
true is a boolean
\end{lstlisting}

\token{println} works the same way from the caller's point of view: it accepts
integers, floats, strings and Booleans alike, and prints each the way its type
requires. How it is written is the subject of Chapter~\ref{chap:templates}.

\subsection{What must differ}

The compiler chooses a function from the arguments of the call, and from nothing
else. Two functions it could not tell apart from their arguments are refused.
The names of the parameters, which the call does not have to write, do not tell
two functions apart.

\begin{lstlisting}[style=coloredverbatimError, escapechar=@, caption=Two functions that differ only by the names of their parameters]
use std::io;

fn area(side: f64)-> f64 {
  side * side
}

fn @\hb{area}@(width: f64)-> f64 { // error
  width * width
}

fn main() {
  println(area(3.0));
}
\end{lstlisting}

The compiler reports \errmsg{E4023}{colliding function definitions \errhole{} and
  \errhole{}}, naming both. The result
type does not tell two functions apart either, since a call does not say which
type it expects.

\begin{lstlisting}[style=coloredverbatimError, escapechar=@, caption=Two functions that differ only by their result]
use std::io;

fn zero()-> i32 {
  0
}

fn @\hb{zero}@()-> f64 { // error
  0.0
}

fn main() {
  println(@\hb{zero()}@); // error
}
\end{lstlisting}

The two declarations collide, as before. The call is reported too: nothing in
\token{zero()} says whether it wants the integer or the floating point zero, and
the compiler says \errmsg{E4280}{\errhole{} called with \{\} works with multiple
  candidates}, the hole being \token{zero}.

\tipbox{Overloading is for functions that do the same job on different inputs,
  so that the reader of a call knows what it does without knowing which of them
  runs. Two functions that do different things deserve different names.}
